\documentclass[12pt]{article}
\usepackage[utf8]{inputenc}
\usepackage[T1]{fontenc}
\usepackage[french]{babel}
\usepackage{amsmath}
\usepackage{amsfonts}
\usepackage{booktabs}
\usepackage{siunitx}


\title{Circuit RL en Régime Sinusoïdal}
\author{Solution détaillée}
\date{}

\begin{document}

\maketitle

\section{Énoncé et Données}
Un circuit série est alimenté par une source de tension sinusoïdale de valeur efficace $V_{eff} = 230 \, \text{V}$ et de fréquence $f = 50 \, \text{Hz}$. Le circuit est composé de :
\begin{itemize}
    \item Une résistance $R = 10 \, \Omega$
    \item Une inductance $L = 100 \, \text{mH} = 0,1 \, \text{H}$
\end{itemize}

\section{Analyse de l'Impédance}
La pulsation du signal est donnée par :
\[ \omega = 2\pi f = 2\pi \cdot 50 = 100\pi \approx 314,16 \, \text{rad/s} \]

L'impédance complexe totale $Z_{eq}$ du circuit série est :
\[ Z_{eq} = R + jL\omega = 10 + j(0,1 \cdot 100\pi) = 10 + j31,4 \, \Omega \]

\textbf{Forme polaire :}
\begin{itemize}
    \item Module : $|Z_{eq}| = \sqrt{R^2 + (L\omega)^2} = \sqrt{10^2 + 31,4^2} \approx 32,95 \, \Omega$
    \item Phase (déphasage) : $\phi = \arctan\left(\frac{L\omega}{R}\right) = \arctan\left(\frac{31,4}{10}\right) \approx 72,3^\circ$
\end{itemize}

\section{Calcul de l'Intensité}
En prenant la tension comme référence de phase ($\underline{V} = 230e^{j0}$), l'intensité complexe est :
\[ \underline{I} = \frac{\underline{V}}{Z_{eq}} = \frac{230}{32,95e^{j72,3^\circ}} \approx 6,98 e^{-j72,3^\circ} \, \text{A} \]
La valeur efficace du courant est donc $I_{eff} \approx 6,98 \, \text{A}$.

\section{Bilan des Puissances}
La puissance complexe $\underline{S}$ est obtenue par le produit de la tension et du conjugué complexe de l'intensité (produit hermitique) :
\[ \underline{S} = \underline{V} \cdot \underline{I}^* = 230 \cdot 6,98e^{+j72,3^\circ} \]
\[ \underline{S} \approx 1605,4e^{j72,3^\circ} \, \text{VA} \]

En utilisant la forme rectangulaire $\underline{S} = P + jQ$ :
\begin{itemize}
    \item \textbf{Puissance active :} $P = \text{Re}(\underline{S}) = 1605,4 \cdot \cos(72,3^\circ) \approx 487 \, \text{W}$
    \item \textbf{Puissance réactive :} $Q = \text{Im}(\underline{S}) = 1605,4 \cdot \sin(72,3^\circ) \approx 1529 \, \text{var}$
\end{itemize}

\section{Synthèse des Résultats}
\begin{table}[h]
    \centering
    \begin{tabular}{lcr}
        \toprule
        Grandeur & Symbole & Valeur \\
        \midrule
        Impédance totale & $|Z_{eq}|$ & $32,95 \, \Omega$ \\
        Courant efficace & $I_{eff}$ & $6,98 \, \text{A}$ \\
        Facteur de puissance & $\cos \phi$ & $0,30$ \\
        Puissance apparente & $S$ & $1605,4 \, \text{VA}$ \\
        \bottomrule
    \end{tabular}
\end{table}

\section{Compensation de la Puissance Réactive}
L'objectif est de relever le facteur de puissance à $\cos \phi' = 0,9$ en installant un condensateur de capacité $C$ en parallèle de l'ensemble $\{R, L\}$.

\subsection{Nouvelle Puissance Réactive Ciblée}
La puissance active $P$ reste inchangée ($P \approx 487 \, \text{W}$). Pour un nouveau $\cos \phi' = 0,9$, l'angle de déphasage devient :
\[ \phi' = \arccos(0,9) \approx 25,84^\circ \]

La nouvelle puissance réactive totale $Q_{tot}$ souhaitée est :
\[ Q_{tot} = P \cdot \tan(\phi') = 487 \cdot \tan(25,84^\circ) \approx 235,8 \, \text{var} \]

\subsection{Calcul de la Capacité $C$}
Le condensateur doit fournir la différence de puissance réactive entre l'état initial ($Q = 1529 \, \text{var}$) et l'état final ($Q_{tot} = 235,8 \, \text{var}$) :
\[ Q_C = Q_{tot} - Q = 235,8 - 1529 \approx -1293,2 \, \text{var} \]

La puissance réactive d'un condensateur est donnée par $Q_C = -C\omega V_{eff}^2$. On en déduit $C$ :
\[ C = \frac{-Q_C}{\omega V_{eff}^2} = \frac{1293,2}{314,16 \cdot 230^2} \]
\[ C \approx 77,8 \cdot 10^{-6} \, \text{F} \approx 77,8 \, \mu\text{F} \]

\subsection{Impact sur le Courant de Ligne}
Grâce à cette compensation, le nouveau courant efficace appelé à la source est :
\[ I'_{eff} = \frac{P}{V_{eff} \cdot \cos \phi'} = \frac{487}{230 \cdot 0,9} \approx 2,35 \, \text{A} \]

\textbf{Conclusion :} En ajoutant ce condensateur, on divise le courant de ligne par près de 3 (passant de $6,98 \, \text{A}$ à $2,35 \, \text{A}$), tout en fournissant la même puissance active à la charge. Cela réduit considérablement les pertes par effet Joule dans les câbles d'alimentation.

\section{Question de Réflexion : Dimensionnement des Conducteurs}

\textbf{Énoncé :} 
Un technicien doit choisir le câble d'alimentation pour le circuit de l'exercice (cas $L = 100 \, \text{mH}$). Il lit sur la plaque signalétique que l'appareil consomme une puissance active $P = 487 \, \text{W}$. 
Pensant bien faire, il calcule le courant en faisant $I = P / V_{eff} = 487 / 230 \approx 2,1 \, \text{A}$ et choisit un câble de très faible section.

\textbf{Question :} 
Expliquez pourquoi ce raisonnement est dangereux et quel est le risque réel pour l'installation.

\textbf{Solution pour l'enseignant :}
Le raisonnement est erroné car le câble est traversé par le courant \textit{réel} (le courant efficace $I_{eff}$), et non par le "courant actif". 
\begin{itemize}
    \item Comme vu précédemment, le courant réel est de $6,98 \, \text{A}$ à cause du faible facteur de puissance ($\cos \phi = 0,3$).
    \item Le courant circulant est donc \textbf{plus de 3 fois supérieur} à ce que le technicien a prévu.
    \item \textbf{Risque :} Le câble va subir un échauffement par effet Joule ($R_{câble} \cdot I^2$) 9 fois supérieur aux prévisions ($3^2 = 9$), ce qui peut entraîner la fonte des isolants et un incendie.
\end{itemize}
\textit{Moralité : On dimensionne toujours une installation en fonction de la puissance apparente $S$ (ou du courant efficace $I_{eff}$), jamais uniquement selon la puissance active $P$.}

\end{document}